% !TeX root = ../深度学习与大语言模型技术原理.tex
% 由主文档通过 \input 载入；请编译主文档。
\setcounter{chapter}{19} % 保留原书章号。
\chapter{提示注入、工具边界与执行隔离}
\label{ch:safety}

\section{威胁来自指令与数据的混淆}

提示注入（prompt injection）试图使模型把不可信内容当成应执行指令。直接注入出现在用户输入，间接注入（indirect prompt injection）可藏在网页、文档、邮件或工具输出中。越狱（jailbreak）通常强调诱导模型绕开行为限制；二者可重叠，但分析重点不同。数据泄露、工具滥用和权限提升是可能结果，不是只有不礼貌文本才构成问题。

首先区分可信指令、用户请求与外部数据。来源标签、引用边界和上下文组织能帮助模型，但纯文本提醒不能构成可靠权限边界。外部网页即使写着“系统消息”，也不应因此取得调用内部工具的能力。安全分析应围绕数据从哪里来、能影响什么动作以及动作由谁授权。

\section{工具调用前后的确定性检查}

调用前应验证工具名、参数模式、资源范围和实际权限。最小权限（least privilege）将只读查询与写操作分开，把数据库、文件和网络权限限制到任务所需范围；业务规则应在程序中执行。例如参数符合 JSON Schema，并不能证明用户有权读取指定记录。

调用后需要限制结果大小、规范化结构、过滤不应返回的敏感字段，并保留来源。结果清洗能减少某些攻击面，却无法保证删除所有恶意语义；因此不可信结果仍保持数据身份。模型生成的后续工具调用应重新检查，不能因“上一步工具成功”而继承无限权限。

应对状态未知、重试和并发写入建立明确规则。审计日志把用户授权、模型建议、校验结果和实际动作串联，便于追查；日志本身也可能含敏感信息，需要控制访问与保留范围。安全与可恢复性在这里共同依赖准确状态记录。

\section{沙箱、资源限制与代码执行}

沙箱（sandbox）为代码、Shell 或浏览器任务提供受限环境。隔离维度包括文件系统、身份、进程、网络、资源和凭证；CPU 时间、内存、输出量与子进程数也要限制。容器有助于封装运行环境，但共享内核等机制意味着“使用了容器”并不自动等于强隔离，应按威胁模型选择合适边界。

不可信代码不应直接获得宿主密钥或无约束网络访问。测试数据与评分程序应避免被候选程序修改；真实业务写入应经过受控接口，而非允许生成代码随意操作生产数据库。清理临时环境后仍需确认外部副作用，因为销毁本地容器不会撤回已经发送的请求。

\section{安全评测与能力保持}

工具护栏（tool guardrail）评测应覆盖越权参数、路径逃逸、恶意检索内容、凭证请求和重复副作用等任务内场景，并记录是否真的执行了动作。只检查最终回答是否拒绝，会漏掉“已执行后再拒绝”的失败。

同时测量正常任务完成率与误拦截率。过度限制可能让系统看似安全却无法完成授权任务，宽松限制则可能降低边界可靠性。把安全结果按攻击入口、目标资源和权限层分类，能比单一总分更准确地定位需要修复的执行环节。

